\chapter{Introduzione}

\section{Contesto e Motivazioni}

Il diametro di un grafo, definito come la massima distanza tra ogni coppia di nodi, 
rappresenta una metrica fondamentale per valutare l'efficienza di una rete.
In ambiti quali le reti di comunicazione, i social network e i sistemi di calcolo multi-core,
mantenere un diametro ridotto è cruciale per minimizzare la latenza e facilitare la diffusione delle informazioni.

Una tecnica comune per migliorare la connettività di una rete esistente consiste nel "shortcutting", 
ovvero nell'aggiunta di un insieme limitato di archi extra, chiamati scorciatoie (\textit{shortcut edges}). 
Il problema classico introdotto in letteratura è il \textit{Bounded Cardinality Minimum Diameter} (BCMD),
il cui obiettivo è aggiungere al più $k$ archi per minimizzare il diametro risultante. 
Tuttavia, le soluzioni ottimali per il BCMD tendono spesso a concentrare i nuovi archi su pochi nodi centrali,
aumentandone drasticamente il grado. 

In molte applicazioni reali, tale aumento è indesiderato poiché può causare sovraccarichi fisici o logistici.
Per ovviare a questo problema, è stata introdotta la variante \textbf{BCMD-$\delta$}, 
che impone un limite superiore $\delta$ all'aumento del grado di ogni singolo vertice.



\section{Definizione Formale del Problema}

In questo lavoro di tesi, estendiamo ulteriormente il problema BCMD-$\delta$ introducendo un vincolo di costo associato agli archi.
Invece di limitare semplicemente il numero di archi aggiunti ($k$), consideriamo una funzione di costo generale e un budget totale.
Chiameremo questo problema \textbf{Budgeted Diameter Minimization with Degree Constraints (BDM-DC)}.
Sia $G = (V, E)$ un grafo non pesato, che può essere orientato o non orientato. Sia $E^c = (V \times V) \setminus E$ 
l'insieme dei potenziali archi scorciatoia aggiungibili al grafo. 
Definiamo i seguenti parametri di input:
\begin{itemize}
    \item Una funzione di costo $c: E^c \to \mathbb{R}^+$ che associa a ogni arco $(u, v) \in E^c$ un valore positivo.
    Nel caso di nodi dotati di coordinate spaziali, $c(u, v)$ corrisponderà tipicamente alla distanza euclidea tra $u$ e $v$.
    \item Un budget totale $B \in \mathbb{R}^+$.
    \item Un limite di grado $\delta \in \mathbb{N}^+$,
    che rappresenta il numero massimo di nuovi archi che possono essere incidenti a ciascun nodo.
\end{itemize}
L'obiettivo è trovare un insieme di archi $M \subseteq E^c$ che minimizzi il diametro del grafo aumentato $G' = (V, E \cup M)$. Formalmente, il problema può essere espresso come:

$$ \min_{M \subseteq E^c} D(G') $$
soggetto ai seguenti vincoli:

\begin{enumerate}
    \item \textbf{Vincolo di Budget:} La somma dei costi degli archi selezionati non deve superare $B$:
    $$ \sum_{e \in M} c(e) \le B $$
    
    \item \textbf{Vincolo di Grado:} Per ogni vertice $v \in V$, l'incremento del suo grado nel grafo $G'$ rispetto al grafo originale $G$ deve essere al più $\delta$:
    $$ deg_{G'}(v) - deg_G(v) \le \delta \quad \forall v \in V $$
\end{enumerate}
Il diametro $D(G')$ è definito come:
$$ D(G') = \max_{u, v \in V} d_{G'}(u, v) $$ 
dove $d_{G'}(u, v)$ è la lunghezza del cammino minimo tra i nodi $u$ e $v$ nel grafo $G'$.

\section{Obiettivi della Tesi}

L'obiettivo principale della ricerca è analizzare la complessità computazionale di questa variante e sviluppare algoritmi (esatti o euristici)
in grado di fornire soluzioni di alta qualità in tempi ragionevoli. In particolare, studieremo come la variazione del budget $B$,
definita proporzionalmente alle caratteristiche del dataset (es. in base al costo medio degli archi), influenzi la capacità di riduzione del 
diametro sotto diversi vincoli di grado $\delta$.